\documentclass[11pt]{hmcpset}
\usepackage[margin=1in]{geometry}
\usepackage{amsmath,amssymb,enumerate,graphicx,geometry}

\newcommand{\ket}[1]{|#1\rangle}
\newcommand{\bra}[1]{\langle #1|}
\newcommand{\braket}[2]{\langle #1|#2\rangle}
\newcommand{\braketop}[3]{\langle #1| #2 | #3 \rangle}
\newcommand{\normsq}[1]{\left|#1\right|^2}
\newcommand{\prob}[2]{\normsq{\braket{#1}{#2}}}
\newcommand{\avg}[1]{\langle #1 \rangle}

\newcommand{\ketvec}[2]{\ket{\!{#1}\!\mathbf{{#2}}}}
\def\bi#1{\textsf{\textit{#1}}}

\name{}
\class{PHYS116}
\assignment{Homework 11}
%\duedate{}

\begin{document}
\problemlist{Townsend: 11.15, 12.2, 12.4, 12.8}


%===========================================================

\begin{problem}[Townsend Problem 11.15 (5\,pts)]

Determine the effect of an external magnetic field on the energy levels of the $n=2$ states of hydrogen when the applied magnetic field $B$ has a magnitude much greater than $10^4$ gauss, in which case the spin-orbit interaction may be neglected as a first approximation. This is the \textbf{Pashen-Bach effect}.
\\

Draw an energy-splitting diagram like Figure 11.10.
\\

\emph{Hint:} Unlike the section problem and most of the section on the Zeeman Effect, you can start with the (11.105) form rather than the (11.106) form of the perturbing Hamiltonian.


\end{problem}
\begin{solution}
	\vfill
\end{solution}
	\pagebreak
	

%===========================================================

\begin{problem}[Townsend Problem 12.2 (5\,pts)]

Two identical, noninteracting spin-$\tfrac{{1}}{{2}}$ particles of mass $ m$ are in the one-dimensional harmonic oscillator for which the Hamiltonian~is
\[
\hat{H} = {\frac{{\hat{p}_{1x}^{2}} }{{2m}}} + {\frac{{1}}{{2}}}m\omega
^{2}\hat{x}_{1}^{2} + {\frac{{\hat{p}_{2x}^{2}} }{{2m}}} +
{\frac{{1}}{{2}}}m\omega ^{2}\hat{x}_{2}^{2}
\]
\begin{enumerate}
	\item[(a)] Determine the ground-state and first-excited-state kets and corresponding
	energies when the two particles are in a total-spin-0 state. What are the
	lowest energy states and corresponding kets for the particles if they are in
	a total-spin-1 state?
	
	\item[(b)] Suppose the two particles interact with a potential energy of interaction
	\[
	V(|x_{1} - x_{2}|) = \begin{cases}
	- V_{0} & | x_{1} - x_{2} | < a \\
	0 & {\rm elsewhere} \\
	\end{cases}
	\]
	Argue what the effect will be on the energies that you determined
	in~($a$), that is, whether the energy of each state moves up,
	moves down, or remains unchanged. \textit{Suggestion}: Examine which spatial wave functions for the total-spin-0 and total-spin-1 states tend to have the particles closer together. Consider, for example, the special case of  $x_1=x_2$.
\end{enumerate}


\end{problem}
\begin{solution}
	\vfill
\end{solution}
	\pagebreak
	

%===========================================================

\begin{problem}[Townsend Problem 12.4 (5\,pts)]

Use the variational principle to estimate the ground-state
energy for the one-dimensional anharmonic oscillator
\[
\hat{H} = {\frac{{\hat{p}_{x}^{2}} }{{2m}}} + b \hat{x}^{4}
\]
Compare your result with the exact result
\[
E_{0} = 1.060b^{1/3} \left( {{\frac{{\hslash ^{2}}}{{2m}}}} \right)^{2/3}
\]


\end{problem}
\begin{solution}
	\vfill
\end{solution}
	\pagebreak
	

%===========================================================

\begin{problem}[Townsend Problem 12.8 (5\,pts)]
	
Use the trial wave function $R(r)=Ne^{-br}$ to estimate the ground-state energy of the three-dimensional isotropic harmonic oscillator, for which $V(r)=\frac{1}{2}\mu \omega^2 r^2$. Compare your estimate with the exact value (see Section~10.5).
	
\end{problem}
\begin{solution}
	\vfill
\end{solution}
\pagebreak


%===========================================================

\end{document}

